% Figure — Générateur de courant suppresseur de variations (a) et générateur auto-polarisé (b) du sujet du 2 avril 2026 : NMOS M1 et M2 sous le seuil, résistance R_Δ (EANA3, chapitre 5).
% D'après le sujet [td04-2026, p. 7-8, figures 1 et 2], redessiné.
% Rendu SVG :
%   latex -output-directory=<dossier temporaire> 05-generateur-courant.tex
%   dvisvgm --no-fonts --exact-bbox -o 05-generateur-courant.svg <dossier temporaire>/05-generateur-courant.dvi
\documentclass[border=4pt]{standalone}
\usepackage[T1]{fontenc}
\usepackage[utf8]{inputenc}
\usepackage{amsmath}
\usepackage{sansmath}
\renewcommand{\familydefault}{\sfdefault}
\usepackage[european, straightvoltages, americancurrents]{circuitikz}
\usetikzlibrary{shapes.geometric}
\begin{document}
\sansmath
\begin{circuitikz}[font=\small, >=latex, fleche/.style={->, thick}]
  % ---------- (a) ----------
  \node at (-1.3,5.6) {(a)};
  \draw (0,0) node[nmos] (M1) {};  \node[right] at (0.45,-0.3) {$M_1$};
  \draw (2.6,1.6) node[nmos] (M2) {};  \node[right] at (3.05,1.3) {$M_2$};
  \draw (M1.D) -- (0,1.6) coordinate (b);
  \draw (b) node[circ]{} to[R, l_=$R_\Delta$, i<_=$I_\Delta$] (0,3.3) coordinate (a);
  \draw (a) node[circ]{} to[R, l=$R_{bias}$, i<=$I_{bias}$] (0,5.3) node[vcc]{$V_{bias}$};
  \draw (a) -- (-1.0,3.3) -- (-1.0,0 |- M1.G) -- (M1.G);
  \draw[fleche] (-0.5,1.75) -- node[left]{$V_\Delta$} (-0.5,3.15);
  \draw (b) -- (M2.G);
  \draw (M2.D) to[short, i<=$I_{out}$] (2.6,3.6) -- (3.6,3.6) node[ocirc]{};
  \draw (M1.S) -- (0,-1.1);  \draw (M2.S) -- (2.6,-1.1);
  \draw (-1.0,-1.1) -- (3.6,-1.1) node[ocirc]{};  \draw (1.3,-1.1) node[ground]{};
  % ---------- (b) ----------
  \begin{scope}[xshift=6.6cm]
    \node at (-1.3,5.6) {(b)};
    \draw (0,4.6) node[pmos, xscale=-1] (Pa) {};  \node[left] at (-0.45,4.6) {$M_{102}$ ($\times 2{,}72$)};
    \draw (2.6,4.6) node[pmos] (Pb) {};  \node[right] at (3.05,4.6) {$M_{101}$};
    \draw (-0.4,5.7) -- (3.0,5.7);  \node[above] at (1.3,5.7) {$V_{DD}$};
    \draw (Pa.S) -- (0,5.7);  \draw (Pb.S) -- (2.6,5.7);
    \draw (Pa.G) -- (Pb.G);
    \draw (1.6,4.6) node[circ]{} -- (1.6,3.4) -- (2.6,3.4) node[circ]{};
    \draw (0,0) node[nmos] (N1) {};  \node[right] at (0.45,-0.3) {$M_1$};
    \draw (2.6,1.6) node[nmos] (N2) {};  \node[right] at (3.05,1.3) {$M_2$};
    \draw (N1.D) -- (0,1.6) coordinate (b);
    \draw (b) node[circ]{} to[R, l_=$R_\Delta$] (0,2.9) coordinate (a);
    \draw (Pa.D) to[short, i=$I_1$] (a);
    \draw (a) node[circ]{} -- (-1.0,2.9) -- (-1.0,0 |- N1.G) -- (N1.G);
    \draw[fleche] (-0.5,1.75) -- node[left]{$V_\Delta$} (-0.5,2.75);
    \draw (b) -- (N2.G);
    \draw (Pb.D) to[short, i=$I_2$] (2.6,3.4) -- (N2.D);
    \draw (N1.S) -- (0,-1.1);  \draw (N2.S) -- (2.6,-1.1);
    \draw (-1.0,-1.1) -- (2.6,-1.1);  \draw (1.3,-1.1) node[ground]{};
  \end{scope}
\end{circuitikz}
\end{document}
